\documentclass[11pt]{article}
\usepackage[a4paper,margin=26mm]{geometry}
\usepackage[T1]{fontenc}
\usepackage{lmodern}
\usepackage{amsmath,amssymb,amsthm}
\usepackage{microtype}
\usepackage{booktabs}
\usepackage{array}
\usepackage{xurl}
\usepackage{hyperref}
\hypersetup{colorlinks=true,linkcolor=blue,urlcolor=blue,citecolor=blue,
 pdftitle={The switching conjecture for main eigenvalues holds for trees of diameter at most 4},
 pdfauthor={Horace Dong}}
\newtheorem{lemma}{Lemma}[section]
\newtheorem{proposition}[lemma]{Proposition}
\newtheorem{corollary}[lemma]{Corollary}
\newtheorem*{thmone}{Theorem 1}
\newtheorem*{conj}{Conjecture}
\theoremstyle{remark}
\newtheorem{remark}[lemma]{Remark}
\newcommand{\Q}{\mathbb Q}
\newcommand{\Z}{\mathbb Z}
\newcommand{\R}{\mathbb R}
\newcommand{\F}{\mathbb F}
\newcommand{\s}{\mathbf s}
\newcommand{\one}{\mathbf 1}
\newcommand{\cL}{\mathcal L}
\newcommand{\cR}{\mathcal R}
\newcommand{\bfa}{\mathbf a}
\newcommand{\NI}{N_{\mathrm I}}
\newcommand{\NII}{N_{\mathrm{II}}}
\newcommand{\Nei}{N_{\mathrm{ei}}}
\newcommand{\sq}{\mathrm{sq}}
\DeclareMathOperator{\rank}{rank}
\DeclareUrlCommand\code{\urlstyle{tt}\def\UrlBreaks{\do\_\do\/\do\.}\def\UrlBigBreaks{}}
\setlength{\emergencystretch}{1.5em}
\title{The switching conjecture for main eigenvalues\\ holds for trees of diameter at most~4}
\author{Horace Dong\\ \small Independent researcher\\
\small\texttt{maxwelldhx@gmail.com}\\
\small ORCID: \href{https://orcid.org/0009-0002-7554-7548}{0009-0002-7554-7548}}
\date{September 2026}
\begin{document}
\maketitle
\begin{abstract}
An eigenvalue of a signed graph is main if its eigenspace is not orthogonal
to the all-ones vector. Akbari, Fran\c{c}a, Ghasemian, Javarsineh and de~Lima
conjectured that every connected graph other than $K_2$ and $K_4-e$ has a
switching, a choice of a sign $\pm1$ at each vertex, for which all eigenvalues
of the resulting signed graph are main, and Akbari, Kumar, Mohar and Pragada
asked for a proof for trees and for connected regular graphs. We prove the
conjecture for every tree of diameter at most~$4$ other than~$K_2$; this class
contains the stars, the double stars and the harmonic trees, for which it was
known. Apart from a few explicit ones, the eigenvalues of such a tree are the
square roots of the roots of a secular equation. If all neighbours of the
centre receive the same sign, each of these whose minimal polynomial is even is
main; along two explicit families of switchings, each
remaining pair of Galois orbits $\pm O$ spoils at most four members, and at
most two when its eigenvalues are irrational. This settles every tree in which a neighbour of the centre has
at least $13$ leaves; the other cases reduce to explicit switchings and to a
finite set of $13{,}376$ trees, which three independently written programs
check. Some abstract steps of the proof are formally verified in Lean~4. The
results were found by AI coding agents working under the author's direction;
their role is described at the end of the paper.
\end{abstract}

\section{Introduction}

\paragraph{Background.}
Let $G$ be a finite simple graph with adjacency matrix $A=A(G)$. A
\emph{switching} of $G$ is a vector $\s\in\{\pm1\}^{V(G)}$. The signed graph
$G^{\s}$ has the edges of $G$, the edge $uv$ carrying the sign $s_us_v$, so
that $A(G^{\s})=D_{\s}AD_{\s}$, where $D_{\s}$ is the diagonal matrix with
diagonal~$\s$~\cite{akmp}. An eigenvalue of a signed graph is \emph{main} if
its eigenspace is not orthogonal to the all-ones vector~$\one$
\cite{cve78,row07}. A regular graph has only one main eigenvalue, whereas almost
all graphs have only main eigenvalues~\cite{ort}. Akbari, Fran\c{c}a,
Ghasemian, Javarsineh and de~Lima~\cite{afgjl} asked whether a suitable
switching always makes all eigenvalues main, and conjectured that it does. We
state the conjecture in the form of~\cite[Conjecture~1.1]{akmp}.

\begin{conj}[{\cite{afgjl}}]
Let $G$ be a connected graph other than $K_2$ and $K_4-e$. Then some switching
$\s$ of $G$ makes every eigenvalue of $G^{\s}$ main.
\end{conj}

Only distinct eigenvalues matter~\cite[Proposition~1.3]{akmp}. The abstract
of~\cite{afgjl} states the conjecture for every graph other than $K_2$ and
$K_4-e$; Akbari, Kumar, Mohar and Pragada~\cite{akmp} state it for connected
graphs and show that connectedness is needed (no switching of $K_2\cup K_1$
works). For trees the two
versions agree.

According to its abstract, \cite{afgjl} proves the conjecture for Cayley,
distance-regular, vertex- and edge-transitive graphs, double stars and paths.
It also covers complete bipartite graphs, hence stars
(\cite[Theorem~4.1]{afgjl}, as cited in~\cite[Section~4]{sy}), and, by a
computer search, all connected graphs on at most nine vertices other than
$K_2$ and $K_4-e$ (as reported in~\cite[Section~1.2]{akmp}
and~\cite[Section~5]{sy}). Shao and
Yuan~\cite{sy} proved it for harmonic trees \cite[Proposition~2.1]{sy},
complete multipartite graphs and the graphs $S_{n,r}$, and
Stani\'c~\cite{sta} gave sufficient conditions for net-regular signed graphs.
The authors of~\cite{akmp} proved two asymptotic versions, with
$n-O(n/(\log n)^{1/4})$ main eigenvalues counted with multiplicity and with
$d-O(d/(\log d)^{1/4})$ distinct main eigenvalues, where $n$ is the order and
$d$ the number of distinct eigenvalues. They single out regular graphs and
trees as unresolved cases and, in~\cite[Problem~1.6]{akmp}, ask to show that
the conjecture ``holds for trees and connected regular graphs of order
$n\ge 3$''.

\paragraph{Result.}
\begin{thmone}
Let $T$ be a tree of diameter at most $4$ other than $K_2$. Then some switching
$\s$ of $T$ makes every eigenvalue of $T^{\s}$ main.
\end{thmone}

The trees of diameter at most $4$ are the trees of radius at most~$2$. They
include the stars, the double stars (all trees of diameter~$3$), the paths
$P_3$, $P_4$, $P_5$, and the harmonic trees of~\cite{sy}, in which a vertex of
degree $a^2-a+1$ has only neighbours of degree $a$ and all other vertices are
leaves. So Theorem~1 contains the known cases above of diameter at most~$4$ and
adds all other trees of diameter~$4$. Since every signature of a tree is a
switching of the all-positive one, Theorem~1 also gives the signature version
of the conjecture~\cite[Conjecture~1.7]{akmp} for these trees.

\paragraph{Method.}
A tree of diameter at most $4$ with at least three vertices is a rooted tree
$T(\bfa)$ of height at most~$2$: a centre $c$ with neighbours $v_1,\dots,v_k$,
where $v_i$ has $a_i$ further neighbours, all of them leaves. Apart from $0$
and the numbers $\pm\sqrt b$ for branch sizes $b\ge1$ that occur more than
once, its eigenvalues are the numbers $\pm\sqrt t$ for the roots $t$ of a
\emph{secular polynomial} $R$ (Lemma~\ref{lem:spec}). Whether an eigenvalue
becomes main after a switching is a simple condition (Lemma~\ref{lem:main}),
which for the secular eigenvalues depends only on the Galois orbit. If all
neighbours of the centre receive the same sign, every such orbit whose minimal
polynomial is even is main (Lemma~\ref{lem:equal}); the other orbits form
\emph{non-even pairs} $\pm O$, with $\theta\in\Q(\theta^2)$ for $\theta\in O$,
which are \emph{integer} pairs if $O\subseteq\Z$ and \emph{irrational} pairs
otherwise. Along two explicit families of switchings, the \emph{rows} of
Section~\ref{sec:rows}, an integer pair spoils at most four members and an
irrational pair at most two, by a linear-independence argument; so a row that
is long compared with the number of non-even pairs contains a good switching
(Corollary~\ref{cor:criteria}). Integer pairs are rare, because they come from
perfect squares that are roots of $R$, and all roots of $R$ but one are smaller
than the largest branch size $b^*=\max_ia_i$ (Lemma~\ref{lem:count}). For
$b^*\ge13$ this makes the row of the largest branches long enough
(Proposition~\ref{prop:large}). For $2\le b^*\le12$ the trees not covered by
the rows lie in an explicit finite set of $13{,}376$ trees, and a computer
search leaves three trees with at most seven vertices, for which we give
switchings (Propositions~\ref{prop:search} and~\ref{prop:three}). Trees with
$b^*\le1$ are treated by explicit switchings (Proposition~\ref{prop:small}).

\paragraph{Scope.}
One step of the proof, Proposition~\ref{prop:search}, is a finite computation,
carried out by three independently written programs (Section~\ref{sec:comp});
all other steps are proved by hand, and some abstract steps are also formally
verified (Section~\ref{sec:lean}). The problem remains open for trees in
general and for regular graphs; our argument relies on the explicit spectrum
of Lemma~\ref{lem:spec}. As supporting evidence, the conclusion of Theorem~1
holds for all $63{,}242{,}254$ trees with $3$ to $24$ vertices, and direct
certificates confirm it for every tree of diameter at most $4$ with at most
$44$ vertices (Section~\ref{sec:comp}).

\paragraph{Prior work.}
The full text of~\cite{afgjl} was not accessible to us: the publisher's page
refused automated access, and we found no other copy. Our account of its
results rests on its abstract, as shown by zbMATH Open, and on the summaries
in~\cite{sy} and~\cite{akmp}. The problem
appeared in version~1 of~\cite{akmp}, submitted to the arXiv on 22 September
2026. We searched for earlier or concurrent work on the conjecture for trees,
in particular for trees of diameter at most~$4$, on 26 September 2026 between
18:48 and 19:49 UTC and between 21:38 and 22:01 UTC, briefly again at 23:00
UTC, and on 27 September 2026 between 00:37 and 00:42 UTC; the last search was
made by the referee agent described in the disclosure. The sources were: the
arXiv (the listing of~\cite{akmp}, which had only version~1, the arXiv API, and
the new submissions in math.CO); zbMATH Open, Semantic Scholar, OpenAlex and
Crossref, including the works citing~\cite{afgjl} and~\cite{sy}; the problem
database MathDB (\url{https://mathdb.com}), whose entry for the conjecture was
still open, with no solutions; MathOverflow and Mathematics Stack Exchange;
GitHub code, commits, issues, pull requests and repositories; public
repositories of AI-assisted work on open problems, among them
\texttt{trureturing}~\cite{tru}; the publication lists of two authors
of~\cite{afgjl}; and web search. We did not search Google Scholar or
MathSciNet, and dblp and some repository search engines could not be queried,
because they returned checks for automated access, which we did not bypass. We
found no proof or claim of the conjecture for all trees of diameter at
most~$4$, or for a class of trees containing them; the tree classes treated in
the work we found are those listed above. According to its abstract, \cite{fbj}
determines the unsigned trees of diameter~$4$ with exactly three main
eigenvalues; it does not concern switchings. A final check on 27 September
2026 between 02:04 and 02:06 UTC (the arXiv, where~\cite{akmp} still had only
version~1; GitHub; MathDB; and MathOverflow) found nothing new. This reports
the coverage of our searches, not a guarantee of priority.

\section{Switchings and main eigenvalues}\label{sec:prelim}
Let $A$ be a real symmetric $n\times n$ matrix. For an eigenvalue $\theta$ of
$A$ let $E_\theta$ be its eigenspace and $P_\theta$ the orthogonal projection
onto~$E_\theta$. For $\s\in\{\pm1\}^n$ the matrix $D_{\s}AD_{\s}$ has the same
eigenvalues, with eigenspaces $D_{\s}E_\theta$ and projections
$D_{\s}P_\theta D_{\s}$.

\begin{lemma}\label{lem:good}
An eigenvalue $\theta$ of $A$ is a main eigenvalue of $D_{\s}AD_{\s}$ if and
only if $P_\theta\s\ne0$.
\end{lemma}
\begin{proof}
The eigenspace $D_{\s}E_\theta$ is orthogonal to $\one$ if and only if the
projection $D_{\s}P_\theta D_{\s}\one=D_{\s}P_\theta\s$ of $\one$ onto it
vanishes.
\end{proof}

We call $\s$ \emph{good} for $A$, or for a graph $G$ with $A=A(G)$, if
$P_\theta\s\ne0$ for every eigenvalue $\theta$ of~$A$. By
Lemma~\ref{lem:good}, Theorem~1 says that every tree of diameter at most $4$
other than $K_2$ has a good switching. If $\s$ is good, so is~$-\s$.

\begin{lemma}[Krylov certificate]\label{lem:krylov}
Let $A$ be a symmetric integer matrix with $d$ distinct eigenvalues, let
$\s\in\Z^n$, and let $K$ be the integer matrix with columns
$\s,A\s,\dots,A^{d-1}\s$. Then $\rank_\Q K$ is the number of eigenvalues
$\theta$ with $P_\theta\s\ne0$. In particular, $\s$ is good if $K$ has rank $d$
over $\F_p$ for some prime~$p$.
\end{lemma}
\begin{proof}
We have $A^j\s=\sum_\theta\theta^jP_\theta\s$, the nonzero vectors
$P_\theta\s$ are orthogonal, and the Vandermonde matrix
$(\theta^j)_{\theta,\,0\le j<d}$ is invertible. So the column space of $K$ is
spanned by the vectors $P_\theta\s$, and its dimension is the number of nonzero
ones. Reduction modulo $p$ does not increase the rank.
\end{proof}

\section{The spectrum of a tree of diameter at most 4}\label{sec:spec}

\paragraph{Notation.}
Let $\bfa=(a_1,\dots,a_k)$, with $k\ge1$, be a multiset of nonnegative
integers. The tree $T(\bfa)$ has a vertex $c$, the \emph{centre}, adjacent to
$v_1,\dots,v_k$, and $v_i$ is adjacent to $a_i$ further vertices, the leaves in
the set~$L_i$. So $T(\bfa)$ has $n=1+k+\sum_ia_i$ vertices. We call $v_i$
together with $L_i$ the $i$th \emph{branch}; it is \emph{bare} if $a_i=0$. Let
$B$ be the set of distinct values of the $a_i$, let $r=|B|$ and
$b^*=\max B$, and for $b\in B$ let $I_b=\{i:a_i=b\}$ and $k_b=|I_b|$; we put
$k_0=0$ if $0\notin B$. We write $B=\{b_1<\dots<b_r\}$.

A tree with $n\ge3$ vertices and diameter at most $4$ has a vertex of
eccentricity at most~$2$; rooting the tree at such a vertex shows that it is
some $T(\bfa)$. Conversely, every $T(\bfa)$ has diameter at most~$4$. The
results of this section and the next hold for every multiset $\bfa$, so the
choice of representation does not matter; a double star, for instance, has
two.

Let
\[
 F(t)=\sum_{b\in B}\frac{k_b}{t-b},\qquad
 R(t)=\prod_{b\in B}(t-b)-\sum_{b\in B}k_b\prod_{b'\in B\setminus\{b\}}(t-b')
 =\prod_{b\in B}(t-b)\,\bigl(1-F(t)\bigr).
\]
Then $R\in\Z[t]$ is monic of degree~$r$; we call it the \emph{secular
polynomial} of~$T(\bfa)$.

\begin{lemma}[spectrum]\label{lem:spec}
Let $T=T(\bfa)$ and $A=A(T)$.
\begin{itemize}
\item[(i)] $R$ has $r$ simple roots $t_1<\dots<t_r$, with $b_j<t_j<b_{j+1}$ for
 $1\le j<r$ and $t_r>b_r$; in particular $t_j>0$ and $t_j\notin B$. For each
 $j$, both $\theta=\sqrt{t_j}$ and $\theta=-\sqrt{t_j}$ are simple eigenvalues
 of~$A$, with eigenvector $x^\theta$ given by $x^\theta_c=1$,
 $x^\theta_{v_i}=\theta/(t_j-a_i)$ and $x^\theta_\ell=1/(t_j-a_i)$ for
 $\ell\in L_i$. We call these $2r$ eigenvalues \emph{secular}.
\item[(ii)] For $b\in B$ with $b\ge1$ and $k_b\ge2$, both $\theta=\pm\sqrt b$ are
 eigenvalues of multiplicity $k_b-1$, and $E_\theta$ consists of the vectors
 with $x_{v_i}=\alpha_i$ and $x_\ell=\alpha_i/\theta$ for $\ell\in L_i$ and
 $i\in I_b$, where $\sum_{i\in I_b}\alpha_i=0$, and with all other
 coordinates~$0$.
\item[(iii)] If $k_0\ge1$, then $\ker A$ consists of the vectors with $x_c=0$,
 with $x_{v_i}=0$ and $\sum_{\ell\in L_i}x_\ell=0$ whenever $a_i\ge1$, and with
 $\sum_{i\in I_0}x_{v_i}=0$. If $k_0=0$, then $\ker A$ consists of the vectors
 with $x_{v_i}=0$ and $\sum_{\ell\in L_i}x_\ell=-x_c$ for every~$i$.
\item[(iv)] $A$ has no other eigenvalues. Hence $A$ has
 $d=2r+2\,|\{b\in B:b\ge1,\ k_b\ge2\}|+[\ker A\ne0]$ distinct eigenvalues.
\end{itemize}
\end{lemma}
\begin{proof}
The equation $Ax=\theta x$ reads
\begin{equation}\label{eq:eigen}
 \theta x_\ell=x_{v_i}\ (\ell\in L_i),\qquad
 \theta x_{v_i}=x_c+\sum_{\ell\in L_i}x_\ell,\qquad
 \theta x_c=\sum_{i=1}^kx_{v_i}.
\end{equation}
Let $\theta\ne0$ and $x\ne0$ be a solution. The first two equations give
$x_\ell=x_{v_i}/\theta$ and $(\theta^2-a_i)x_{v_i}=\theta x_c$. If $x_c\ne0$ we
may take $x_c=1$; then $\theta^2\ne a_i$ for all $i$,
$x_{v_i}=\theta/(\theta^2-a_i)$, and the third equation becomes
$F(\theta^2)=1$, that is, $R(\theta^2)=0$. Conversely, these formulas solve
\eqref{eq:eigen} whenever $R(\theta^2)=0$ and $\theta^2\notin B$. If $x_c=0$,
then $x_{v_i}=0$ unless $a_i=\theta^2$, so $\theta^2=b\in B$ with $b\ge1$, and
the third equation gives $\sum_{i\in I_b}x_{v_i}=0$; this is~(ii), and it
needs $k_b\ge2$. For $\theta=0$ the equations say $x_{v_i}=0$ when $a_i\ge1$,
$x_c+\sum_{\ell\in L_i}x_\ell=0$ for every $i$, and $\sum_ix_{v_i}=0$; a bare
branch forces $x_c=0$, and (iii) follows.

The value $R(b_j)=-k_{b_j}\prod_{i\ne j}(b_j-b_i)$ has sign $(-1)^{r-j+1}$, and
$R(t)\to+\infty$ as $t\to\infty$. So $R$ changes sign on each interval
$(b_j,b_{j+1})$ and on $(b_r,\infty)$, which gives $r$ roots as stated; there
are no others, since $\deg R=r$. As no $t_j$ lies in $B$, the eigenvalues in
(i), (ii) and (iii) are distinct, and those in (i) are simple. Finally the
dimensions add up to~$n$: they are $2r$ in~(i), $2\sum_{b\ge1}(k_b-1)$ in~(ii),
and $\sum_{a_i\ge1}(a_i-1)+k_0-1$ if $k_0\ge1$, or $1+\sum_i(a_i-1)$ if
$k_0=0$, in~(iii); in both cases the total is $1+k+\sum_ia_i$.
\end{proof}

\paragraph{The effect of a switching.}
For a switching $\s$ of $T(\bfa)$ we write $s_c=s(c)$, $\sigma_i=s(v_i)$ and
$\lambda_i=\sum_{\ell\in L_i}s(\ell)$, the \emph{leaf sum} of branch~$i$
($\lambda_i=0$ for a bare branch), and for $b\in B$
\[
 S_b=\sum_{i\in I_b}\sigma_i,\qquad \Lambda_b=\sum_{i\in I_b}\lambda_i,\qquad
 A_b=s_ck_b+\Lambda_b .
\]

\begin{lemma}[main eigenvalues after switching]\label{lem:main}
A switching $\s$ of $T=T(\bfa)$ is good if and only if the following three
conditions hold.
\begin{itemize}
\item[(S)] For every secular eigenvalue $\theta$, with $t=\theta^2$,
 \[
 G_{\s}(\theta):=\sum_{b\in B}\frac{A_b+S_b\theta}{t-b}\ne0 .
 \]
\item[(L)] For every $b\in B$ with $b\ge1$ and $k_b\ge2$ and each
 $\theta\in\{\sqrt b,-\sqrt b\}$, the numbers $\sigma_i+\lambda_i/\theta$
 with $i\in I_b$ are not all equal.
\item[(Z)] If $0$ is an eigenvalue: when $k_0\ge1$, some branch with
 $a_i\ge2$ has leaves of both signs, or the $\sigma_i$ with $i\in I_0$ are not
 all equal; when $k_0=0$, some branch has leaves of both signs, or
 $s_c\ne\sum_i\varepsilon_i$, where $\varepsilon_i$ is the common sign of the
 leaves in~$L_i$.
\end{itemize}
\end{lemma}
\begin{proof}
By Lemma~\ref{lem:spec}, $\s$ is good if and only if $\s$ is not orthogonal to
$x^\theta$ for each secular $\theta$, nor to $E_{\pm\sqrt b}$, nor to
$\ker A$. First,
$\s\cdot x^\theta=s_c+\sum_i(\sigma_i\theta+\lambda_i)/(t-a_i)$; writing
$s_c=s_cF(t)$ and grouping the branches by size gives $G_{\s}(\theta)$.
Second, for $x\in E_\theta$ as in Lemma~\ref{lem:spec}(ii) we have
$\s\cdot x=\sum_{i\in I_b}\alpha_i(\sigma_i+\lambda_i/\theta)$, which vanishes
for all $\alpha$ with $\sum\alpha_i=0$ exactly when the coefficients are all
equal. Third, let $k_0\ge1$. By Lemma~\ref{lem:spec}(iii), $\s\perp\ker A$ if
and only if $\s$ is constant on each $L_i$ with $a_i\ge2$ and on
$\{v_i:i\in I_0\}$. Let $k_0=0$. If two leaves $\ell,\ell'$ of one branch have
opposite signs, then $e_\ell-e_{\ell'}\in\ker A$ and
$\s\cdot(e_\ell-e_{\ell'})\ne0$. Otherwise $\s\cdot x=x_c(s_c-\sum_i\varepsilon_i)$
for $x\in\ker A$, and $\ker A$ contains a vector with $x_c=1$.
\end{proof}

\paragraph{Orbits.}
The polynomial $R(x^2)\in\Z[x]$ is monic and even, and by
Lemma~\ref{lem:spec}(i) its roots are the $2r$ secular eigenvalues, which are
distinct. Let $f$ be a monic irreducible factor of $R(x^2)$ in $\Q[x]$; we call
the set of its roots an \emph{orbit}. We say that $\s$ \emph{fails at} an orbit
$O$ if $G_{\s}(\theta)=0$ for some $\theta\in O$. A rational number that is a
root of $R$ is an integer, since $R$ is monic with integer coefficients.

\begin{lemma}[orbits]\label{lem:orbits}
Let $f$ be as above, $\theta$ a root of $f$, and $t=\theta^2$.
\begin{itemize}
\item[(a)] If $f(-x)=f(x)$, then $\theta\notin\Q(t)$. We call the orbit of
 $f$ \emph{even}. (As $0$ is not a root of $R(x^2)$, the case $f(-x)=-f(x)$
 does not occur.)
\item[(b)] If $f(-x)\ne f(x)$, then $\theta\in\Q(t)$, and $\pm f(-x)$ is another
 irreducible factor of $R(x^2)$, whose orbit is the negative of the orbit of~$f$. We call
 these two orbits a \emph{non-even pair}. If $\deg f=1$, then $\theta\in\Z$ and
 $t$ is a perfect square that is a root of $R$ (an \emph{integer pair});
 otherwise $t\notin\Q$ (an \emph{irrational pair}).
\item[(c)] If an integer root $t$ of $R$ is not a perfect square, then
 $\{\sqrt t,-\sqrt t\}$ is an even orbit.
\item[(d)] Let $\NI$ and $\NII$ be the numbers of integer and of irrational
 pairs, and $\Nei$ the number of roots of $R$ that are even integers but not
 perfect squares. Then $\NI+2\NII+\Nei\le r$.
\item[(e)] For every switching $\s$, if $\s$ fails at an orbit $O$, then
 $G_{\s}$ vanishes on all of~$O$.
\end{itemize}
\end{lemma}
\begin{proof}
Let $h\in\Q[t]$ be the minimal polynomial of $t$. Then $h\mid R$ and
$f\mid h(x^2)$. If $\theta\notin\Q(t)$, then
$\deg f=[\Q(\theta):\Q]=2\deg h$, so $f=h(x^2)$ is even. If $\theta\in\Q(t)$,
then $\deg f=\deg h$; if $f$ were even, say $f=f_1(x^2)$, then $f_1(t)=0$ would
give $\deg f=2\deg f_1\ge2\deg h$, which is impossible. So $f$ is not even, and
$\pm f(-x)$, which divides $R(x^2)$ because $R(x^2)$ is even, is a different
irreducible factor. If $\deg f=1$, then $\theta$ is a rational root of the
monic integer polynomial $R(x^2)$, hence an integer. If $\deg f\ge2$, then
$\deg h=\deg f\ge2$ and $t\notin\Q$. This proves (a) and~(b), and (c) holds
because $x^2-t$ is irreducible. For~(d), the minimal polynomials of the roots
of $R$ are distinct irreducible factors of~$R$: each integer pair uses a factor
$t-q^2$, each irrational pair a factor of degree at least~$2$, and each root
counted by $\Nei$ a factor $t-m$ with $m$ not a square; their degrees add up to
at most $\deg R=r$. For~(e), write $G_{\s}(x)=g(x)/\prod_{b\in B}(x^2-b)$ with
$g\in\Q[x]$; the denominator does not vanish at secular eigenvalues, and if
$g(\theta)=0$ for one root $\theta$ of $f$, then $f\mid g$.
\end{proof}

The following reduction is the starting point of the proof.

\begin{lemma}[equal signs next to the centre]\label{lem:equal}
Suppose that $\sigma_1=\dots=\sigma_k=\sigma$. Then
$G_{\s}(\theta)=\sigma\theta+P_{\s}(t)$ with
$P_{\s}(t)=\sum_{b\in B}A_b/(t-b)\in\Q(t)$ for every secular $\theta$. Hence
$\s$ fails at no even orbit, and $\s$ fails at a non-even pair, that is, at one
of its two orbits, if and only if $P_{\s}(t)\in\{\theta,-\theta\}$, where
$\theta$ is any root of the pair and $t=\theta^2$.
\end{lemma}
\begin{proof}
Here $S_b=\sigma k_b$, so $\sum_bS_b\theta/(t-b)=\sigma\theta F(t)=\sigma\theta$.
In an even orbit $\theta\notin\Q(t)$, so $\sigma\theta+P_{\s}(t)\ne0$. The last
statement holds because $G_{\s}(\theta)=0$ means $P_{\s}(t)=-\sigma\theta$ and
$G_{\s}(-\theta)=0$ means $P_{\s}(t)=\sigma\theta$.
\end{proof}

\section{Two rows of switchings}\label{sec:rows}

In this section we assume $b^*\ge2$, that is, some branch has at least two
leaves.

\paragraph{Standard leaf sums.}
For every branch $i$ with $a_i\ge1$ we fix a number $\lambda^\circ_i$ with
$|\lambda^\circ_i|\le a_i$ and $\lambda^\circ_i\equiv a_i\pmod 2$, group by
group, after ordering each $I_b$:
\begin{itemize}
\item if $b=1$: $\lambda^\circ_i=-1$, except $\lambda^\circ_i=1$ for the
 smallest $i\in I_1$ when $k_1\ge2$;
\item if $b\ge3$ is odd: $\lambda^\circ_i=-1$, except that the two smallest
 $i\in I_b$ get $1$ and $-3$ when $k_b\ge2$;
\item if $b\ge2$ is even: $\lambda^\circ_i=0,-2,0,-2,\dots$ along~$I_b$.
\end{itemize}
Then (P1) for $b\ge1$ with $k_b\ge2$ the $\lambda^\circ_i$ with $i\in I_b$ are
not all equal, and (P2) for $b\ge2$ the smallest $i\in I_b$ has
$|\lambda^\circ_i|<b$, so the leaves of that branch receive both signs whenever
their sum is~$\lambda^\circ_i$. Put $\Lambda^\circ_b=\sum_{i\in
I_b}\lambda^\circ_i$ for $b\ge1$ and
\[
 U(t)=\sum_{b\in B,\,b\ge1}\frac{\Lambda^\circ_b}{t-b},\qquad
 U_\beta(t)=U(t)-\frac{\Lambda^\circ_\beta}{t-\beta}.
\]

\begin{lemma}[realizable leaf sums]\label{lem:real}
Let $\beta\ge1$ and $k\ge1$, and let $\cL_\beta(k)$ be the set of integers
$\Lambda$ for which there are $\lambda_1,\dots,\lambda_k\in\{-\beta,-\beta+2,
\dots,\beta\}$ with $\sum_i\lambda_i=\Lambda$, not all equal if $k\ge2$, and,
if $\beta\ge2$, with $|\lambda_i|<\beta$ for some~$i$. Then for $\beta\ge2$,
$\cL_\beta(k)$ consists of the $\Lambda\equiv k\beta\pmod2$ with
$|\Lambda|\le k\beta-2$, except $\Lambda=0$ when $\beta=k=2$; so
$|\cL_\beta(k)|=k\beta-1-[\beta=k=2]$. For $\beta=1$, $\cL_1(k)$ consists of the
$\Lambda\equiv k\pmod 2$ with $|\Lambda|\le k-2$ if $k\ge2$, and
$\cL_1(1)=\{\pm1\}$; so $|\cL_1(k)|=k-1$ for $k\ge2$ and $|\cL_1(1)|=2$.
\end{lemma}
\begin{proof}
Every realizable $\Lambda$ satisfies the stated parity and $|\Lambda|\le
k\beta$, and $|\Lambda|=k\beta$ would force all $\lambda_i$ to equal $\pm\beta$.
For $(\beta,k,\Lambda)=(2,2,0)$ the only candidates are $(0,0)$ and
$(2,-2)$, which are excluded. The case $\beta=1$ is clear. Conversely, let
$\beta\ge2$ and let $\Lambda$ be as stated. Put $M=(k\beta-\Lambda)/2$, the
number of leaves of sign $-1$; then $1\le M\le k\beta-1$. For $k=1$ take
$\lambda_1=\Lambda$. For $k\ge2$ distribute $M$ as evenly as possible, $m_i\in
\{\lfloor M/k\rfloor,\lceil M/k\rceil\}$, and put $\lambda_i=\beta-2m_i$. If
$k\nmid M$, the $m_i$ are not all equal and some $m_i$ lies in $[1,\beta-1]$.
If $M=kq$, then $1\le q\le\beta-1$, and we use $(q+1,q-1,q,\dots,q)$ instead;
some $m_i$ lies in $[1,\beta-1]$ unless $k=2$, $q=1$ and $\beta=2$, which is
the excluded case.
\end{proof}

We write $\cL_\beta=\cL_\beta(k_\beta)$ for $\beta\in B$ with $\beta\ge1$.

\begin{lemma}[leaf row]\label{lem:leafrow}
Let $\beta\in B$ with $\beta\ge1$. For $\varepsilon\in\{\pm1\}$ and
$\Lambda\in\cL_\beta$ let $\s_{\varepsilon,\Lambda}$ be a switching with
$s_c=\varepsilon$, with $\sigma_i=1$ for every branch (bare branches included),
with leaf sums $\lambda_i=\lambda^\circ_i$ for $a_i\ge1$, $i\notin I_\beta$,
and with leaf sums on $I_\beta$ as in Lemma~\ref{lem:real} with total
$\Lambda$ (branch $i$ gets $(a_i-\lambda_i)/2$ leaves of sign $-1$). These
$2|\cL_\beta|$ switchings form the \emph{leaf row} of~$\beta$. Every member
satisfies (L) and~(Z), and for every secular $\theta$, with $t=\theta^2$,
\begin{equation}\label{eq:leafG}
 G_{\s_{\varepsilon,\Lambda}}(\theta)=\theta+\varepsilon+U_\beta(t)
 +\frac{\Lambda}{t-\beta}.
\end{equation}
No member fails at an even orbit. An integer pair makes at most four members
fail, at most two for each~$\varepsilon$, and an irrational pair makes at most
two members fail.
\end{lemma}
\begin{proof}
(L) holds by (P1) for $b\ne\beta$ and by Lemma~\ref{lem:real} for $b=\beta$,
since all $\sigma_i=1$. For (Z): if $\beta\ge2$, Lemma~\ref{lem:real} gives a
branch of size $\beta$ with $|\lambda_i|<\beta$, whose leaves have both signs;
if $\beta=1$, the group $b^*\ge2$ keeps its standard leaf sums, and (P2) gives
such a branch. Since $A_b=\varepsilon k_b+\Lambda_b$ with $\Lambda_0=0$,
Lemma~\ref{lem:equal} gives $G=\theta+\varepsilon F(t)+\sum_{b\ge1}\Lambda_b/
(t-b)$, which is~\eqref{eq:leafG}, and no member fails at an even orbit. By
Lemma~\ref{lem:equal}, $\s_{\varepsilon,\Lambda}$ fails at a non-even pair if
and only if $\varepsilon+U_\beta(t)+\Lambda w\in\{\theta,-\theta\}$, where
$w=1/(t-\beta)\ne0$; for fixed $\varepsilon$ the left-hand side is injective in
$\Lambda$, so at most two $\Lambda$ fail. If the pair is irrational, then
$t\notin\Q$, so $w\notin\Q$ and $1,w$ are linearly independent over~$\Q$. The
member fails at the orbit of $\theta$ exactly when
$\varepsilon+\Lambda w=-\theta-U_\beta(t)$, and at the orbit of $-\theta$
exactly when $\varepsilon+\Lambda w=\theta-U_\beta(t)$; each of these equations
has at most one solution $(\varepsilon,\Lambda)\in\Z^2$.
\end{proof}

\begin{lemma}[bare row]\label{lem:barerow}
Suppose $k_0\ge1$. For $\varepsilon\in\{\pm1\}$ and $0\le m\le k_0$ let
$\s_{\varepsilon,m}$ be a switching with $s_c=\varepsilon$, with leaf sums
$\lambda^\circ_i$ on every branch with $a_i\ge1$, with $\sigma_i=1$ on these
branches, and with $\sigma_i=-1$ on exactly $m$ of the bare branches and
$\sigma_i=1$ on the others. These $2(k_0+1)$ switchings form the \emph{bare
row}. Every member satisfies (L) and~(Z), and for every secular $\theta$
\begin{equation}\label{eq:bareG}
 G_{\s_{\varepsilon,m}}(\theta)=\varepsilon+U(t)+\theta-\frac{2m}{\theta}.
\end{equation}
An even orbit makes at most one member fail, and only if its $t=\theta^2$ is
an even integer. An integer pair makes at most four members fail, and an
irrational pair at most two.
\end{lemma}
\begin{proof}
(L) holds by (P1), and (Z) by (P2) applied to the group~$b^*$. Now $S_0=k_0-2m$
and $S_b=k_b$ for $b\ge1$, so $\sum_bS_b/(t-b)=F(t)-2m/t=1-2m/t$, and
\eqref{eq:bareG} follows as before. In an even orbit, $1$ and $\theta$ are
linearly independent over $\Q(t)$, so $G=(\varepsilon+U(t))+\theta(1-2m/t)$
vanishes only if $t=2m$ and $\varepsilon=-U(t)$; this determines the member.
Let $(\theta,-\theta)$ be a non-even pair. The member fails at the orbit of
$\theta$ if and only if $\varepsilon-2m/\theta=-\theta-U(t)$, and at the orbit of
$-\theta$ if and only if $\varepsilon+2m/\theta=\theta-U(t)$. For fixed
$\varepsilon$ each equation has at most one solution~$m$, so at most four
members fail. If the pair is irrational, then $1/\theta\notin\Q$, so each
equation has at most one solution $(\varepsilon,m)\in\Z^2$.
\end{proof}

\begin{corollary}\label{cor:criteria}
Let $b^*\ge2$. Then $T(\bfa)$ has a good switching if
\begin{itemize}
\item[(a$'$)] $|\cL_\beta|\ge2\NI+\NII+1$ for some $\beta\in B$ with
 $\beta\ge1$, or
\item[(b$'$)] $k_0\ge1$ and $2(k_0+1)>4\NI+2\NII+\Nei$.
\end{itemize}
In case (a$'$) some member of the leaf row of $\beta$ is good; in case (b$'$)
some member of the bare row.
\end{corollary}
\begin{proof}
By Lemma~\ref{lem:leafrow}, at most $4\NI+2\NII$ of the $2|\cL_\beta|$ members
of the leaf row fail at some orbit, so under (a$'$) some member fails at none.
It satisfies (S) because every secular eigenvalue lies in an orbit, and (L) and
(Z) by Lemma~\ref{lem:leafrow}; so it is good by Lemma~\ref{lem:main}. Under
(b$'$) the same argument applies to the bare row, with at most
$4\NI+2\NII+\Nei$ failing members by Lemma~\ref{lem:barerow} and
Lemma~\ref{lem:orbits}(c).
\end{proof}

The primes match the labels printed by our programs, where (a) and (b)
denote an earlier count with four failing members for every non-even pair.

\section{Proof of Theorem 1}\label{sec:proof}

\subsection{No branch with two leaves}

\begin{proposition}\label{prop:small}
Let $T=T(\bfa)$ have $n\ge3$ vertices and $b^*\le1$, so that
$T=T(1^{k_1},0^{k_0})$: the centre has $k_0$ leaf neighbours and $k_1$
neighbours with one further leaf each. Then $T$ has a good switching, namely:
\begin{itemize}
\item[(a)] if $k_1=0$ (the star $K_{1,k_0}$, $k_0\ge2$): $s_c=1$, and $-1$ on one
 leaf (on two leaves if $k_0=4$) and $+1$ on the other leaves;
\item[(b)] if $k_1=1$ and $k_0\ge2$: $s_c=1$; $-1$ on one leaf neighbour of $c$
 and $+1$ on the others; $+1$ on the other neighbour $v$ of $c$ and $-1$ on the
 leaf of~$v$;
\item[(c)] if $k_1\ge2$ and $k_0\ge2$: $s_c=1$; $-1$ on one leaf neighbour of
 $c$ and $+1$ on the others; $+1$ on the $k_1$ other neighbours of $c$; $-1$ on
 one of their leaves and $+1$ on the others;
\item[(d)] if $k_0=1$: one of the two switchings $\s^{\pm}$ with $s_c=\pm1$,
 all $\sigma_i=1$ and leaf sums $\lambda^\circ_i$ (defined as in
Section~\ref{sec:rows});
\item[(e)] if $k_0=0$ (a spider with legs of length~$2$): the switching~$\s^+$.
\end{itemize}
\end{proposition}
\begin{proof}
In all cases we use Lemma~\ref{lem:main}.

(a) Here $R(t)=t-k_0$, so $t=k_0$, and $G=(k_0+E\theta)/t=1+E/\theta$, where
$E\in\{k_0-2,0\}$ is the sum of the signs of the leaves. $G=0$ would need
$E^2=k_0$, and $(k_0-2)^2=k_0$ only for $k_0\in\{1,4\}$, while $E=0$ for
$k_0=4$. (L) is void, and (Z)
holds because the leaf signs are not all equal.

(b) Here $R(t)=t^2-(k_0+2)t+k_0$ and
$G=1+(k_0-2)/\theta+(\theta-1)/(t-1)$. Substituting $k_0=t(t-2)/(t-1)$, which
is $R(t)=0$, gives
$G=(\theta^2-2)(\theta^2+\theta-1)/(\theta(\theta^2-1))$. Now $t=2$ is not a
root, as $R(2)=-k_0$. If $\theta^2+\theta-1=0$, then $t=\theta^2=1-\theta$
satisfies $t^2-3t+1=0$, and subtracting this from $R(t)=0$ gives
$(k_0-1)(1-t)=0$, which is impossible because $k_0\ge2$ and $t\ne1$. (L) is void, and (Z) holds as in~(a).

(c) Here $R(t)=h(t):=t^2-(1+k_0+k_1)t+k_0$. With $A_0=k_0$, $S_0=k_0-2$,
$A_1=2k_1-2$ and $S_1=k_1$, and using $F(t)=1$, we get $G=P+\theta Q$ with
$P=(t+k_1-3)/(t-1)$ and $Q=(t-2)/t$. If $G$ vanishes at $\theta$ or at
$-\theta$, then $P^2=tQ^2$, that is,
$f(t):=t(t+k_1-3)^2-(t-2)^2(t-1)^2=0$. Since $h(0)=k_0>0$, $h(1)=-k_1<0$,
$h(k_0+k_1)=-k_1<0$ and $h(k_0+k_1+1)=k_0>0$, the roots of $h$ lie in $(0,1)$
and in $(k_0+k_1,k_0+k_1+1)$; so $h$ has no integer root and is irreducible
over~$\Q$. Hence $h\mid f$, that is, the remainder $c_1t+c_0$ of $f$ modulo
$h$ vanishes, where
\begin{align*}
 c_1&=-k_0^3-3k_0^2k_1+6k_0^2-3k_0k_1^2+12k_0k_1-13k_0-k_1^3+7k_1^2-12k_1+8,\\
 c_0&=k_0^3+2k_0^2k_1-6k_0^2+k_0k_1^2-7k_0k_1+13k_0-4 .
\end{align*}
As polynomials in $k_0$, both have constant leading coefficients $\mp1$, so a
common root forces their resultant to vanish:
\[
 \operatorname{Res}_{k_0}(c_1,c_0)=(k_1-1)\bigl(k_1^6-11k_1^5+60k_1^4-148k_1^3
 +224k_1^2-192k_1+64\bigr)=0 .
\]
For $k_1\ge2$ the first factor is nonzero, and the sextic has no integer root
(an integer root would divide $64$, and none of $\pm1,\pm2,\dots,\pm64$ is a
root). So (S) holds. (L) holds because the leaves of the $k_1$ branches do not
all have the same sign, and (Z) holds as in~(a).

(d) Here $R(t)=t^2-(k_1+2)t+1$ has discriminant $k_1(k_1+4)$, strictly between
$(k_1+1)^2$ and $(k_1+2)^2$, so $R$ is irreducible, $\NI=0$, and there is at most
one non-even pair. By Lemma~\ref{lem:equal}, $\s^{\pm}$ fails at no even orbit,
and $P_{\s^+}(t)=1/t+2[k_1\ge2]/(t-1)$ and $P_{\s^-}(t)=P_{\s^+}(t)-2F(t)=
P_{\s^+}(t)-2$. If both failed at the pair, then $P_{\s^+}(t)$ and
$P_{\s^+}(t)-2$ would both lie in $\{\theta,-\theta\}$, so $2\in\{\pm2\theta\}$
and $t=1$, which is not a root. (L) holds by (P1). The eigenvalue $0$ does not
occur, by Lemma~\ref{lem:spec}(iii), since $k_0=1$ and all $a_i\le1$.

(e) Here $R(t)=t-k_1-1$ and $P_{\s^+}(t)=2[k_1\ge2]/(t-1)$. For $k_1=1$,
$G=\theta\ne0$. For $k_1\ge2$, $G(\mu)=\mu+2/(\mu^2-1)$ for each secular
eigenvalue $\mu$, and $G(\mu)=0$ means $\mu^3-\mu+2=0$. This covers both
$\mu=\theta$ and $\mu=-\theta$, and both are eigenvalues. The polynomial
$x^3-x+2$ has no rational root, so it is irreducible, and its discriminant is
$-104<0$, so it has only one real root. But every conjugate of an eigenvalue of
a symmetric rational matrix is again an eigenvalue, hence real; so no
eigenvalue satisfies $\mu^3-\mu+2=0$. (This step is formalised; see
Section~\ref{sec:lean}. Directly: $\mu^2=k_1+1$ would give $\mu k_1=-2$, so
$\mu^2=4/k_1^2<k_1+1$.) (L) holds by (P1). For (Z), each leaf group is a
single leaf, and $\sum_i\varepsilon_i=\Lambda^\circ_1\in\{-1,2-k_1\}$ differs
from $s_c=1$.
\end{proof}

\subsection{Counting non-even pairs}

\begin{lemma}\label{lem:count}
Let $b^*\ge1$, let $g=|\{0,1,\dots,b^*\}\setminus B|$, so that $r=b^*+1-g$, and
let $\sq$ be the number of integers $q\ge1$ with $q^2\le b^*-1$ and
$q^2\notin B$. Then $\sq\le\lfloor\sqrt{b^*-1}\rfloor$, $\sq\le g$,
$\NI\le\sq+1$ and $\NI+2\NII+\Nei\le r$.
\end{lemma}
\begin{proof}
By Lemma~\ref{lem:spec}(i), the roots $t_1,\dots,t_{r-1}$ lie in
$(b_1,b^*)\subseteq[0,b^*)$, and no root of $R$ lies in~$B$. An integer pair has
$t=q^2$ with $q\ge1$ (Lemma~\ref{lem:orbits}(b)); if $t\ne t_r$, then
$q^2\le b^*-1$ and $q^2\notin B$. Hence $\NI\le\sq+1$. The squares counted by
$\sq$ are distinct elements of $\{1,\dots,b^*-1\}\setminus B$, so $\sq\le g$;
the bound $\sq\le\lfloor\sqrt{b^*-1}\rfloor$ is clear, and the last inequality
is Lemma~\ref{lem:orbits}(d).
\end{proof}

\subsection{Large branches}

\begin{proposition}\label{prop:large}
If $b^*\ge13$, then (a$'$) holds with $\beta=b^*$. Hence some member of the
leaf row of $b^*$ is good.
\end{proposition}
\begin{proof}
With the notation of Lemma~\ref{lem:count},
\[
 2(2\NI+\NII)=3\NI+(\NI+2\NII)\le3(\sq+1)+(b^*+1-g)\le2\sq+b^*+4
 \le2\lfloor\sqrt{b^*-1}\rfloor+b^*+4 .
\]
For $b^*\ge13$ we have $2\lfloor\sqrt{b^*-1}\rfloor+7\le b^*$: if $b^*\le16$,
then $\lfloor\sqrt{b^*-1}\rfloor=3$; if $b^*\ge17$, then
$q=\lfloor\sqrt{b^*-1}\rfloor\ge4$ satisfies $q^2+1\le b^*$ and
$q^2+1-(2q+7)=(q-1)^2-7>0$. So $2(2\NI+\NII)\le2b^*-3$, and since the
left-hand side is even, $2\NI+\NII\le b^*-2$. By Lemma~\ref{lem:real},
$|\cL_{b^*}|=k_{b^*}b^*-1\ge b^*-1\ge2\NI+\NII+1$, which is (a$'$). The good
member of the leaf row is given by Corollary~\ref{cor:criteria}, that is, by
the counts of Lemma~\ref{lem:leafrow}.
\end{proof}

For $b^*=12$ the same estimate gives only $2\NI+\NII\le11=b^*-1$; for this
reason the smaller values of $b^*$ need a search.

\subsection{Branches with two to twelve leaves}

For a set $B\subseteq\{0,1,\dots,b^*\}$ containing $b^*$, with $r=|B|$ and
$\sq=\sq(B)$ as in Lemma~\ref{lem:count}, put
\[
 \nu(B)=\min(\sq+1,\,r),\qquad
 M(B)=2\nu(B)+\bigl\lfloor(r-\nu(B))/2\bigr\rfloor .
\]

\begin{lemma}\label{lem:region}
Let $2\le b^*\le12$, and suppose that $T(\bfa)$ satisfies neither (a$'$), for
any $\beta$, nor (b$'$). Then $|\cL_\beta|\le M(B)$ for every $\beta\in B$ with
$\beta\ge1$, and $2(k_0+1)\le3\nu(B)+r$ if $k_0\ge1$. Let $\cR$ be the set of
multisets $\bfa$ with $2\le b^*\le12$ that satisfy these two conditions. Then
$\cR$ is finite; it has $13{,}376$ elements, and each of them has $n\le124$.
\end{lemma}
\begin{proof}
By Lemma~\ref{lem:count}, $\NI\le\nu(B)$ and $\NII\le\lfloor(r-\NI)/2\rfloor$.
The function $x\mapsto2x+\lfloor(r-x)/2\rfloor$ is increasing, so
$2\NI+\NII\le M(B)$, and failing (a$'$) gives $|\cL_\beta|\le2\NI+\NII\le M(B)$.
Failing (b$'$) gives $2(k_0+1)\le4\NI+2\NII+\Nei=3\NI+(\NI+2\NII+\Nei)\le
3\nu(B)+r$. There are finitely many sets $B$, and for each of them the first
condition bounds every multiplicity $k_b$ with $b\ge1$, because
$|\cL_b(k)|\ge k-1$ by Lemma~\ref{lem:real}, and the second bounds~$k_0$.
The count of $\cR$ and the bound on $n$ were obtained by enumeration
(Section~\ref{sec:comp}).
\end{proof}

\begin{proposition}\label{prop:search}
Every multiset in $\cR$ except $(2,2)$, $(2,0,0)$ and $(3)$ satisfies (a$'$)
or~(b$'$).
\end{proposition}

This is the computer-assisted step of the proof. It was carried out by three
independently written programs, which evaluate $\NI$, $\NII$ and $\Nei$
exactly; see Section~\ref{sec:comp}.

\begin{proposition}\label{prop:three}
The following switchings are good; the three trees have $7$, $6$ and $5$
vertices.
\begin{itemize}
\item[(a)] $T(2,2)$: $s_c=\sigma_1=\sigma_2=1$; the leaves of $v_1$ get $+1,+1$
 and those of $v_2$ get $-1,-1$.
\item[(b)] $T(2,0,0)=D(2,2)$, the double star with two leaves at each of its
 two central vertices: $s_c=1$; $+1$ on the two leaf neighbours of $c$; $-1$ on the
 third neighbour $v$ of $c$, and $+1,-1$ on the two leaves of~$v$.
\item[(c)] $T(3)=K_{1,4}$, rooted at a leaf $c$: $s_c=1$; $+1$ on the centre $v$
 of the star; $+1,-1,-1$ on the other three leaves.
\end{itemize}
\end{proposition}
\begin{proof}
(a) $R(t)=t-4$, so the secular eigenvalues are $\pm2$, and
$G=(2+2\theta)/2=1+\theta\in\{3,-1\}$. For (L), at $\theta=\pm\sqrt2$ the
numbers $1+2/\theta$ and $1-2/\theta$ differ. For (Z), $k_0=0$ and the leaf
signs are constant on each branch, with $\sum_i\varepsilon_i=0\ne1=s_c$.
(b) $R(t)=(t-1)(t-4)$, and
$G=(2+2\theta)/t+(1-\theta)/(t-2)$ takes the values $4$, $-2$, $1$, $1$ at
$\theta=1,-1,2,-2$. (L) is void, and (Z) holds because the leaves of $v$ have
both signs. (c) $R(t)=t-4$ and $G=\theta/(t-3)=\theta\ne0$; (L) is void, and (Z)
holds because the three leaves adjacent to $v$ have both signs.
\end{proof}

These three trees are also covered by known results: $K_{1,4}$ is a star,
$T(2,0,0)$ a double star, and $T(2,2)$ has seven vertices.

\subsection{End of the proof}

\begin{proof}[Proof of Theorem~1]
The tree $K_1$ is trivial, since the switching $(1)$ is good, and $K_2$ is
excluded. Let $T$ have $n\ge3$ vertices; then $T=T(\bfa)$ for some $\bfa$
(Section~\ref{sec:spec}). If $b^*\le1$, Proposition~\ref{prop:small} applies. If $b^*\ge13$,
Proposition~\ref{prop:large} applies. If $2\le b^*\le12$ and (a$'$) or (b$'$)
holds, Corollary~\ref{cor:criteria} applies. Otherwise $\bfa\in\cR$ by
Lemma~\ref{lem:region}, so $\bfa$ is $(2,2)$, $(2,0,0)$ or $(3)$ by
Proposition~\ref{prop:search}, and Proposition~\ref{prop:three} applies.
\end{proof}



\section{Computations}\label{sec:comp}
The programs and logs are in version~1.7.0 of~\cite{repo}, directory
\code{certificates/switching/}; those of the referee agent (see the
disclosure) begin with \code{rv_}. They are written in Python with
SymPy~\cite{sympy}, except \code{swtrees.c} and \code{swtrees_print.c} (in C),
and all ran on one computer (macOS, arm64).

\paragraph{The finite search.}
Proposition~\ref{prop:search} was checked by three separately written
programs.
\begin{itemize}
\item[(1)] \code{final_region.py} enumerates $\cR$ set by set, as in
 Lemma~\ref{lem:region}. For each tree it finds the integer roots of $R$
 exactly, by evaluating $R$ at every integer in the intervals of
 Lemma~\ref{lem:spec}(i), the last one being $(b^*,b^*+k]$; this gives $\NI$,
 $\Nei$ and the number $N_{\mathrm{rat}}$ of rational roots. It then uses
 $\NII\le\lfloor(r-N_{\mathrm{rat}})/2\rfloor$, which is valid because an
 irrational pair uses two irrational roots and because (a$'$) and (b$'$) only
 become harder as $\NII$ grows. All but ten trees pass (a$'$) or (b$'$) this
 way; for these ten, $R(x^2)$ is factored over~$\Q$, which gives the exact
 $\NII$, and the three trees of Proposition~\ref{prop:three} remain. The run
 takes half a second.
\item[(2)] \code{final_region_indep.py} enumerates a cruder superset of
 $7{,}892{,}382$ multisets, with only the bounds
 $2\NI+\NII\le\lfloor(2\lfloor\sqrt{b^*-1}\rfloor+b^*+4)/2\rfloor$ and
 $2(k_0+1)\le3\lfloor\sqrt{b^*-1}\rfloor+b^*+4$, which follow from
 Lemma~\ref{lem:count} as in the proofs of Proposition~\ref{prop:large} and
 Lemma~\ref{lem:region}, and no pruning by~$B$. It finds the integer
 roots of $R$ by a floating-point screen of $F(t)-1$ at every integer $t$,
 confirming each hit in exact rational arithmetic; a true root gives
 $|F(t)-1|$ far below the threshold. For the trees that fail this cheap test it
 bounds $\NII$ from above: it factors $R(t)$ over $\Q$ and counts an orbit as
 even only if its own Rabin test shows that $h(x^2)$ is irreducible modulo a
 prime not dividing the leading coefficient. It finds the same ten trees and
 the same three survivors.
\item[(3)] The referee's \code{rv_region.py} shares no code with (1) and (2).
 It enumerates $\cR$ in another way, over all multiplicity vectors under global
 caps, with the bounds of Lemma~\ref{lem:region} computed by brute force over
 the feasible $(\NI,\NII,\Nei)$, and writes the set to
 \code{rv_region_list.txt}; a comparison made separately from the program,
 and repeated when the certificates were packaged, shows that this set equals
 that of~(1). It computes $\NI$,
 $\NII$ and $\Nei$ exactly for every tree and certifies each orbit by its own
 tests: irreducibility by a Rabin test, evenness by a simple root of $h$ modulo
 a prime that is a quadratic non-residue, and non-evenness by an explicit
 factorisation $h(x^2)=\pm g(x)g(-x)$. Over $\cR$, $13{,}199$ trees have no
 non-even pair, $176$ have one and one tree has two. It finds the same three
 survivors.
\end{itemize}
Independently of Sections~\ref{sec:spec}--\ref{sec:proof}, the referee's
\code{rv_certify.py} found for each of the $13{,}376$ trees of $\cR$ a switching
that is good by Lemma~\ref{lem:krylov}, with ranks modulo $2^{31}-1$ and
$2147483629$ and with $d$ computed exactly, in agreement with
Lemma~\ref{lem:spec}(iv); for the $3{,}029$ trees with $n\le40$ it also computed
the rank over~$\Q$. Our \code{final_region_witness.py} checked that for each of
the other $13{,}373$ trees of $\cR$ the row given by (a$'$) or (b$'$) contains
a certified good member and has no more failing members than
Lemmas~\ref{lem:leafrow} and~\ref{lem:barerow} allow.

\paragraph{Explicit switchings.}
\code{five_refined.py} checked the switchings of Proposition~\ref{prop:three}
by the rank over~$\Q$ and, independently, by testing $(\mu/p)(A)\s\ne0$ for
every irreducible factor $p$ of the minimal polynomial $\mu$ of~$A$, which for
rational $\s$ is equivalent to goodness. The referee's
\code{rv_exceptional.py} tested all $2^n$ switchings of the three trees: $72$ of
$128$, $16$ of $64$ and $12$ of $32$ are good, including ours. Its
\code{rv_families.py} redid the computation in
Proposition~\ref{prop:small}(c) and certified the switchings of
Proposition~\ref{prop:small} for $2\le k_0\le40$ and $k_1\le40$, for stars with
$k_0\le80$, and for $k_0\in\{0,1\}$ with $k_1\le80$.

\paragraph{Checks of the lemmas.}
These exact checks are not needed for the proof. \code{test_refined_rows.py}
tested all $41{,}039$ leaf and bare rows ($380{,}520$ switchings) of the
$11{,}535$ multisets with $5\le n\le27$ and $b^*\ge2$ against
Lemmas~\ref{lem:leafrow} and~\ref{lem:barerow}; \code{test_bounds.py} confirmed
Lemma~\ref{lem:count}, (a$'$) at $\beta=b^*$ when $b^*\ge13$, and that only the
trees of Proposition~\ref{prop:three} fail (a$'$) and (b$'$), for the
$375{,}820$ multisets with $4\le n\le44$ and $b^*\ge2$. The referee checked
Lemmas~\ref{lem:spec} and~\ref{lem:main} on $400$ trees against exact ranks
(\code{rv_spectrum.py}), and the rows of $69$ trees with non-even pairs and of
the $9{,}114$ multisets with $b^*\ge2$ and $n\le26$ (\code{rv_rows.py},
\code{rv_endtoend.py}). The bounds of Lemma~\ref{lem:leafrow} are attained, for
instance by $T(15)$ (four failing members) and by $T(18,0,0)$ (two, from one
irrational pair).

\paragraph{Supporting evidence.}
Two programs checked the trees with $3\le n\le24$: both of them for
$n\le22$, and implementation~B alone for $n=23$ and~$24$. Each of these
$63{,}242{,}254$ trees has a good switching, certified by
Lemma~\ref{lem:krylov} with $d$ computed exactly, and the counts for each $n$
agree with sequence A000055 of~\cite{oeis}. Implementation~B,
\code{swtrees.c}, generates the trees with the algorithm of Wright, Richmond,
Odlyzko and McKay~\cite{wrom} (a port of the NetworkX~\cite{nx} generator),
computes $d$ exactly from $\gcd(\varphi,\varphi')$ (found modulo a $61$-bit
prime and confirmed by exact division over~$\Z$) and computes ranks modulo
$2^{31}-1$. Implementation~A, \code{swtrees_A.py}, is independent in every
component (canonical augmentation, SymPy for $d$, exact rank over $\Q$ by
fraction-free elimination); it covers $n\le21$, and $n=22$ with trees from the
NetworkX generator (\code{swtrees_A_nx.py}). Separately,
\code{verify_diam4.py} certified a good switching for each of the $376{,}324$
multisets with $3\le n\le44$, which include every tree of diameter at most $4$
with at most $44$ vertices, with $p=2^{61}-1$ and $d$ from
Lemma~\ref{lem:spec}; for $n\le30$ it also computed $d$ from
$\gcd(\varphi,\varphi')$, with no mismatch. The referee's \code{rv_certify.py}
confirmed the $23{,}023$ multisets with $n\le30$ by the rank over~$\Q$.
Non-even pairs are rare: of the $376{,}324$ multisets, $373{,}837$ have none,
$2{,}442$ have one, $44$ have two and one has three
(\code{noneven_census.py}).

\section{Formal verification}\label{sec:lean}
Some abstract steps of the proof are formalised in Lean~4 (version 4.33.1)
with Mathlib~\cite{mathlib} (commit \texttt{0df444a3}), in the files
\code{SwitchingThmH.lean}, \code{SwitchingRowCount.lean},
\code{SwitchingDiam4Cubic.lean} and \code{SwitchingWalkProfile.lean} of
directory \code{lean/Research/} in version~1.7.0 of~\cite{repo}. Each has an
audit file (\code{...Audit.lean}) that restates or prints the main statements;
those of the last three files also check small instances. Table~\ref{tab:lean} lists the main
declarations. There $K$ is any field of characteristic~$0$, and ``$1$ and $w$
are linearly independent over~$\Q$'' is stated as: $a+cw=0$ with $a,c\in\Q$
implies $a=c=0$.

\begin{table}[ht]\centering\small
\begin{tabular}{@{}>{\raggedright\arraybackslash}p{0.3\textwidth}
 >{\raggedright\arraybackslash}p{0.42\textwidth}
 >{\raggedright\arraybackslash}p{0.2\textwidth}@{}}\toprule
Declaration & Statement & Used in\\\midrule
\code{SwitchingThmH.two_sqrt_add_seven_le} & for $b\in\mathbb N$: $b\ge13
 \Rightarrow2\lfloor\sqrt{b-1}\rfloor+7\le b$ &
 Proposition~\ref{prop:large}\\\addlinespace
\code{SwitchingThmH.row_budget} & for $b,\sq,g,N_1,N_2\in\mathbb N$: if
 $b\ge13$, $\sq\le\lfloor\sqrt{b-1}\rfloor$, $\sq\le g$, $N_1\le\sq+1$ and
 $N_1+2N_2+g\le b+1$, then $2N_1+N_2+1\le b-1$ &
 Proposition~\ref{prop:large}\\\addlinespace
\code{SwitchingThmH.sol_subsingleton} & $w,u\in K$, $1$ and $w$ linearly
 independent over $\Q$ $\Rightarrow$ at most one $(\varepsilon,\Lambda)\in\Z^2$
 has $\varepsilon+\Lambda w=u$ & Lemmas~\ref{lem:leafrow},
 \ref{lem:barerow}, irrational pairs\\\addlinespace
\code{SwitchingThmH.indep_of_irrational} & $w\in K$ and $w\ne q$ for all
 $q\in\Q$ $\Rightarrow$ $1$ and $w$ linearly independent over $\Q$ & the
 same\\\addlinespace
\code{SwitchingRow.card_bad_le_two} & $u,w,\theta\in K$, $w\ne0$ $\Rightarrow$
 at most two $\Lambda\in\Z$ have $u+\Lambda w\in\{\theta,-\theta\}$ (a
 bound on \code{Set.ncard}; the set is finite) &
 Lemmas~\ref{lem:leafrow}, \ref{lem:barerow}, integer pairs\\\addlinespace
\code{SwitchingDiam4.no_eigenvalue_root_cubic} & $A\in\Q^{n\times n}$,
 $A^{\mathsf T}=A$, $\theta\in\R$ a root of the characteristic polynomial of~$A$
 $\Rightarrow\theta^3-\theta+2\ne0$ & Proposition~\ref{prop:small}(e), for
 $\theta$ and~$-\theta$\\\addlinespace
\code{SwitchingWalkProfile.main_after_switching} & Remark~\ref{rem:regular};
 also \code{annihilator}, \code{eigen_nonorth},
 \code{walkRegular_main_after_switching} & not used for Theorem~1\\\bottomrule
\end{tabular}
\caption{The main Lean declarations.}\label{tab:lean}
\end{table}

The formal statements are abstract. \code{row_budget} is the arithmetic of
Proposition~\ref{prop:large}; its hypotheses come from Lemma~\ref{lem:count}
and are proved on paper. Of the counts in Lemmas~\ref{lem:leafrow}
and~\ref{lem:barerow}, only two abstract steps are formalised:
\code{card_bad_le_two}, used for a fixed $\varepsilon$ after renaming $u$, $w$
and the target pair, and the linear-independence step
\code{sol_subsingleton} with \code{indep_of_irrational}. That a member fails
at a pair exactly when one of the displayed equations holds, and that the
relevant $w$, namely $1/(t-\beta)$ or $\mp2/\theta$, is irrational for an
irrational pair, is proved on paper only. In particular the refined count
itself, that an irrational pair spoils at most two members of a row, is not
formalised as a statement about trees. \code{no_eigenvalue_root_cubic} is
exactly the algebraic step of Proposition~\ref{prop:small}(e), applied to
$\theta$ and to~$-\theta$. (\code{SwitchingRowCount.lean} also proves
\code{card_bad_le_one_of_ne}, which we do not use.) Everything else is proved
on paper only, or by computer (Section~\ref{sec:comp}); in particular
Lemmas~\ref{lem:good}, \ref{lem:krylov}, \ref{lem:spec}--\ref{lem:equal},
\ref{lem:real}, \ref{lem:count} and~\ref{lem:region},
Corollary~\ref{cor:criteria}, the rest of Propositions~\ref{prop:small}
and~\ref{prop:large}, the finite search of Proposition~\ref{prop:search} and
the checks of Proposition~\ref{prop:three} are not formalised.

The recorded axiom audits (\texttt{\#print axioms}) report only
\texttt{propext}, \texttt{Classical.choice} and \texttt{Quot.sound} for every
declaration in Table~\ref{tab:lean}; the files contain no \texttt{sorry},
\texttt{admit} or \texttt{native\_decide} and declare no axioms. Each file and its
audit file were first rebuilt by \texttt{lake build}, with exit code~$0$, in a
fresh directory that contained only the project configuration, a root file and
these two sources, the prebuilt Mathlib packages being linked in; this was done
on 26 September 2026 at 20:12, 22:31, 23:28 and 23:53 UTC for
\code{SwitchingWalkProfile}, \code{SwitchingDiam4Cubic},
\code{SwitchingRowCount} and \code{SwitchingThmH}, and these logs are in
\code{certificates/switching/}. After the last of these rebuilds, one line of
the module docstring of \code{SwitchingThmH.lean} was reworded; no Lean code
changed. The files as released are covered by the clean replay of the whole
repository at version~1.7.0, made on 27 September 2026 between 01:32 and 01:47
UTC in a fresh directory, with the sources and the configuration copied from
the release tree and the Mathlib packages of the pinned commit taken from a
shared cache (logs \code{logs/clean-replay-v1.7.0-*} of~\cite{repo}):
\texttt{lake build} of all modules completed its $8758$ jobs with exit
code~$0$; the eleven audit files of the repository report no axioms other than
the three above; \texttt{leanchecker} checked the four switching modules, with
exit code~$0$; and the SHA-256 checksums of all $66$ source and configuration
files, listed in \code{logs/lean-sources-sha256.txt} and including that of the
released \code{SwitchingThmH.lean}, were confirmed in the replay directory.

\section{Further remarks and questions}\label{sec:rem}

\begin{remark}[how many vertices must be switched?]\label{rem:an}
For a tree $T$ with $n\ge3$ vertices let $\tau(T)$ be the least value of
$\min(|U|,n-|U|)$ over the sets $U$ of vertices for which the switching that
is $-1$ exactly on $U$ is good, and let $a(n)$ be the maximum of $\tau(T)$ over
the trees with $n$ vertices. The computations of Section~\ref{sec:comp} give,
for $n=3,4,\dots,24$,
\[
 a(n)=1,1,2,2,2,2,2,2,2,2,3,2,3,3,3,3,3,3,3,3,3,4 .
\]
The lower bounds come from explicit trees, checked exactly by the rank over
$\Q$ and by the test $(\mu/p)(A)\s\ne0$ (\code{lower_bounds.py},
\code{crosscheck_galois.py}). The upper bounds rest on implementation~B,
confirmed by implementation~A for $n\le20$ (\code{swtrees_A_tiers.py}). At
$n=24$, B found $39$ trees without a certified switching of at most three
vertices; exact checks (\code{swtrees_print.c}, \code{verify_hard24.py}) show
that none has one and that each has a good switching with $|U|=4$; the checks
of seven of these trees were completed when the certificates were packaged
(\code{certificates/switching/checks-addition/}). For
instance, the tree on the vertices $0,\dots,12$ with edges $0$--$1$, $1$--$2$,
$2$--$3$, $3$--$4$, $0$--$5$, $5$--$6$, $6$--$7$, $5$--$8$, $0$--$9$,
$9$--$10$, $9$--$11$ and $0$--$12$ needs three switched vertices. So a proof
for all trees cannot switch at most three vertices. We do not know whether
$a(n)$ is bounded.
\end{remark}

\begin{remark}[pairing in bipartite graphs]\label{rem:G}
Lemma~\ref{lem:equal} reflects a general fact. Let $G$ be bipartite, with
$A=\begin{pmatrix}0&C\\C^{\mathsf T}&0\end{pmatrix}$, let $\theta\ne0$ be an
eigenvalue whose minimal polynomial is even, and let $\s$ be rational. Then
$P_\theta\s=0$ if and only if $P_{-\theta}\s=0$. Indeed, $E_{\pm\theta}$
consists of the vectors $(u,\pm C^{\mathsf T}u/\theta)$ with
$u\in\ker(CC^{\mathsf T}-\theta^2)$, a space with a basis over $\Q(\theta^2)$,
and $\s\cdot(u,\pm C^{\mathsf T}u/\theta)=a\pm\theta b$ with
$a,b\in\Q(\theta^2)$, while $\theta\notin\Q(\theta^2)$. So in a tree only the
eigenvalues with $\theta\in\Q(\theta^2)$, such as the roots of the factors
$x\pm1$ and $x^2\pm x-1$ of the characteristic polynomial of the tree in
Remark~\ref{rem:an}, impose separate conditions for $\theta$ and~$-\theta$.
\end{remark}

\begin{remark}[regular graphs]\label{rem:regular}
We do not treat regular graphs, but record a sufficient condition that is
formalised in \code{SwitchingWalkProfile.lean}. Let $A$ be real symmetric of
order $n\ne2$ with constant row sums~$k$, and suppose that
$(A^j)_{vv}=\sum_uw_u(A^j)_{uu}$ for all $j\ge0$, for an index $v$ and weights
$w_u>0$. Then $\s=\one-2e_v$ is good for~$A$: if $f(A)\s=0$ for a polynomial $f$,
then $f(k)\one=2Me_v$ for $M=f(A)$, summing coordinates gives $(n-2)f(k)=0$, so
$Me_v=0$, and $0=(M^2)_{vv}=\sum_uw_u\|Me_u\|^2$ forces $M=0$; now take
$f=\mu_A/(x-\theta)$. For walk-regular graphs ($w_u=1/n$) this already follows
from Stani\'c's criterion \cite[Corollary~4.6]{sta}, since then
$(P_\theta)_{vv}=m_\theta/n>0$, where $m_\theta$ is the multiplicity. We do not
claim this case as new.
\end{remark}

\paragraph{Questions.}
The problem of~\cite{akmp} remains open for trees in general, beginning with
the trees of diameter $5$ and~$6$, and so does the question whether the number
$a(n)$ of Remark~\ref{rem:an} is bounded.

\medskip
\noindent Thanks are due to Saieed Akbari, Hitesh Kumar, Bojan Mohar and
Shivaramakrishna Pragada for the problem studied here, and to the authors
of~\cite{afgjl} for the conjecture.

\section*{Tool and computational resource disclosure}
This work was carried out on 26 and 27 September 2026 (UTC) in a workflow
directed by the author, using AI coding agents: Anthropic Claude Code, with the
model Claude Opus~5.5. Claude Code agents selected the problem, found the
theorem and its proof, wrote the programs and the Lean code, carried out the
literature searches described in the introduction, and drafted this text.
Throughout, ``we'' refers to this workflow; the reading, searching and
computing described in the paper were done by the agents. An adversarial
checking agent (Claude Code, the same model), launched by the agent that found
the proof, re-derived Lemmas~\ref{lem:leafrow}, \ref{lem:barerow}
and~\ref{lem:region} and Proposition~\ref{prop:large} and repeated the finite
search; it found one sign error in a displayed equation, which did not affect
the count and has been corrected. Its programs were not kept. An independent
referee agent (Claude Code, the same model), which had not produced the work,
then refereed all proofs and the correspondence between the Lean statements and
the text, recomputed the finite search, the exceptional trees and further
checks with its own programs (Section~\ref{sec:comp}), and repeated the
literature search. Its five minor corrections, concerning a cross-reference in
the case $b^*\ge13$, the description of the Lean coverage, a shorter proof of
Proposition~\ref{prop:small}(c), the sign $-\theta$ in
Proposition~\ref{prop:small}(e) and the description of the independent
computations, have been applied. A second independent agent (Claude Code, the
same model), which had produced neither the work nor the referee's report,
then reviewed this paper, re-derived its proofs and recomputed the
computations of the proof with its own programs; its minor corrections,
concerning the description of the Lean verification and some wordings, have
been applied. The programs, logs and Lean code cited in this paper are in
version~1.7.0 of the repository~\cite{repo} (directories
\code{certificates/switching/}, \code{lean/Research/} and \code{logs/}). The
author is not a
professional mathematician.
The author has read the paper.
The correctness of Theorem~1 rests on the written proofs and on the finite
computation of Proposition~\ref{prop:search}, on which three independently
written programs agree; the Lean formalisation covers only the steps listed in
Section~\ref{sec:lean}. AI systems are not authors of this paper; the author
takes full responsibility for its content.

{\footnotesize
\begin{thebibliography}{99}\setlength{\itemsep}{0pt}
\bibitem{afgjl}
S.~Akbari, F.~A.~M. Fran\c{c}a, E.~Ghasemian, M.~Javarsineh and L.~S. de~Lima,
The main eigenvalues of signed graphs, \emph{Linear Algebra Appl.}
\textbf{614} (2021), 270--280, \url{https://doi.org/10.1016/j.laa.2020.04.029}.
\bibitem{akmp}
S.~Akbari, H.~Kumar, B.~Mohar and S.~Pragada, The switching conjecture for
main eigenvalues is asymptotically true, arXiv:2609.27046v1, 22 September
2026.
\bibitem{cve78}
D.~M. Cvetkovi\'c, The main part of the spectrum, divisors and switching of
graphs, \emph{Publ.\ Inst.\ Math.\ (Beograd) (N.S.)} \textbf{23(37)} (1978),
31--38.
\bibitem{repo}
H.~Dong, \emph{ai4math-results}, repository,
\url{https://github.com/Horace-Maxwell/ai4math-results}, version~1.7.0
(2026); archived at Zenodo under the concept DOI
\url{https://doi.org/10.5281/zenodo.22970254}.
\bibitem{fbj}
F.~A.~M. Fran\c{c}a, A.~E. Brondani and D.~A. Jaume, Graphs with exactly three
main eigenvalues, \emph{Discrete Appl.\ Math.} \textbf{378} (2026), 49--64,
\url{https://doi.org/10.1016/j.dam.2025.07.003}.
\bibitem{nx}
A.~A. Hagberg, D.~A. Schult and P.~J. Swart, Exploring network structure,
dynamics, and function using NetworkX, in \emph{Proceedings of the 7th Python in
Science Conference} (SciPy 2008), 2008, 11--15.
\bibitem{mathlib}
The mathlib Community, The Lean mathematical library, in \emph{Proc.\ 9th ACM
SIGPLAN Int.\ Conf.\ Certified Programs and Proofs} (CPP '20), ACM, 2020,
\url{https://doi.org/10.1145/3372885.3373824}.
\bibitem{sympy}
A.~Meurer et al., SymPy: symbolic computing in Python, \emph{PeerJ Comput.\
Sci.} \textbf{3} (2017), e103, \url{https://doi.org/10.7717/peerj-cs.103}.
\bibitem{oeis}
OEIS Foundation Inc., \emph{The On-Line Encyclopedia of Integer Sequences},
\url{https://oeis.org}, sequence A000055, accessed 26 September 2026.
\bibitem{ort}
S.~O'Rourke and B.~Touri, On a conjecture of Godsil concerning controllable
random graphs, \emph{SIAM J.\ Control Optim.} \textbf{54} (2016), 3347--3378.
\bibitem{row07}
P.~Rowlinson, The main eigenvalues of a graph: a survey, \emph{Appl.\ Anal.\
Discrete Math.} \textbf{1} (2007), 445--471.
\bibitem{sy}
Z.~Shao and X.~Yuan, Some signed graphs whose eigenvalues are main, \emph{Appl.\
Math.\ Comput.} \textbf{423} (2022), Paper No.~127014,
\url{https://doi.org/10.1016/j.amc.2022.127014}; also arXiv:2106.07878v3,
whose numbering we use.
\bibitem{sta}
Z.~Stani\'c, Main eigenvalues of real symmetric matrices with application to
signed graphs, \emph{Czechoslovak Math.\ J.} \textbf{70(145)} (2020),
1091--1102, \url{https://doi.org/10.21136/CMJ.2020.0147-19}.
\bibitem{tru}
The Omega Institute, \texttt{trureturing} repository,
\url{https://github.com/the-omega-institute/trureturing}, accessed 26 and 27
September 2026.
\bibitem{wrom}
R.~A. Wright, B.~Richmond, A.~Odlyzko and B.~D. McKay, Constant time generation
of free trees, \emph{SIAM J.\ Comput.} \textbf{15} (1986), 540--548.
\end{thebibliography}}
\end{document}
